% paper/section_4_temporal_spatial.tex
% Drop-in replacement for §4 of paper/skeleton.tex.
% This is the strongest novel theoretical contribution per the research summary.

\section{The Temporal--Spatial Defense-in-Depth Distinction}
\label{sec:temporal-spatial}

Existing agent guardrails — LlamaFirewall \cite{saxe2025llamafirewall},
AgentSpec \cite{wang2025agentspec}, AGrail \cite{luo2025agrail},
ShieldAgent \cite{chen2025shieldagent} — share a defense topology we call
\emph{spatial}: independent enforcement points operate in parallel along the
request path at a fixed time-step. Each point can independently veto. The
guarantee is that, at the moment of decision, the attacker must defeat
multiple disjoint detectors.

We argue that spatial defense alone is insufficient for systems that admit
runtime extension — plugins, MCP servers, dynamically registered tools — and
that an additional axis we call \emph{temporal} is required. This section
formalizes the distinction, gives a necessity claim with proof sketch, and
states Kavach's coverage theorem.

\subsection{Definitions}

\begin{definition}[Enforcement point]
\label{def:enforcement-point}
An \emph{enforcement point} \(E\) is a pair \((P, A)\) where \(P\) is a
predicate over an event \(e\) (e.g., a tool-call request) and \(A\) is an
action that conditionally modifies the system's behavior in response to
\(P(e)\). The action may be: ALLOW (no modification), BLOCK (the event is
not delivered to its handler), MODIFY (the event payload is rewritten), or
ESCALATE (defer to a human). An enforcement point is \emph{authoritative}
during a time interval \(I\) if no event in \(I\) reaches its handler
without first being evaluated by \(E\).
\end{definition}

\begin{definition}[Spatial defense-in-depth]
\label{def:spatial}
A guardrail \(G\) is \emph{spatially layered} with degree \(k\) if, for
every event \(e\), \(G\) evaluates \(e\) against \(k\) independent
enforcement points \(\{E_1, \dots, E_k\}\) at the same time-step, and
delivers \(e\) only if all \(k\) return ALLOW. The guardrail is sound under
the failure of any \(k-1\) detectors.
\end{definition}

LlamaFirewall is a canonical spatial defense with \(k=3\) (PromptGuard 2 +
AlignmentCheck + CodeShield). AgentSpec is spatial with \(k\) equal to the
number of declared rules. ShieldAgent is spatial with \(k\) equal to the
number of clauses in its rule circuit.

\begin{definition}[Temporal defense-in-depth]
\label{def:temporal}
A guardrail \(G\) is \emph{temporally layered} with respect to an event
class \(\mathcal{E}\) if, for every event \(e \in \mathcal{E}\) and every
component \(C\) that may emit or handle \(e\), every enforcement point
\(E_i \in G\) is authoritative on the time interval during which \(C\)
exists. Equivalently: \(G\) cannot lose authority through the registration,
loading, or lifecycle of a component that emits or handles \(e\).
\end{definition}

The two axes are orthogonal: a guardrail can be spatial without being
temporal (typical case), temporal without being spatial (a single,
unbreakable enforcement point), or both. Kavach is both.

\subsection{The necessity of temporal layering}

\begin{claim}[Spatial defense is insufficient for runtime-extensible systems]
\label{claim:necessity}
Let \(S\) be a system that admits runtime registration of components that
emit events in some class \(\mathcal{E}\), and let \(G\) be a spatially
layered guardrail of any degree \(k\). Then there exists an attacker
strategy that causes \(S\) to handle some event \(e \in \mathcal{E}\)
without \(G\) evaluating \(e\), unless \(G\) is also temporally layered
with respect to \(\mathcal{E}\).
\end{claim}

\paragraph{Proof sketch.} Spatial layering guarantees that, at any
\emph{single} time-step \(t\), if \(G\)'s enforcement points are all
authoritative at \(t\), then any event evaluated at \(t\) is checked
against all \(k\) points. The guarantee says nothing about whether
\(G\)'s enforcement points \emph{remain} authoritative at \(t' > t\).

Suppose at \(t_0\), \(G\) is initialized and its enforcement points are
authoritative. Suppose at \(t_1 > t_0\), an attacker registers a component
\(C\) into \(S\). \(C\) may emit events in \(\mathcal{E}\) at any
\(t_2 > t_1\). For \(G\) to still evaluate those events, \(G\)'s
enforcement-point authority must extend to \(t_2\) — that is, \(G\) must
be temporally layered with respect to \(\mathcal{E}\). If \(G\) is only
spatially layered, the registration of \(C\) is outside \(G\)'s defense
contract.

A concrete instantiation arises in OpenClaw issue \#5513: the platform's
typed-hook runner snapshots the plugin registry at initialization, and any
plugin (including a hypothetical security plugin) that registers after the
snapshot has its hooks silently no-op. A spatial guardrail registered after
initialization is \emph{visible} in the registry but not \emph{authoritative}
at the time of subsequent tool calls. The bug's existence is empirical
evidence for Claim~\ref{claim:necessity}: spatial defense with one rail
that is non-authoritative is no defense at all. \(\square\)

\subsection{Kavach's temporal coverage}

Kavach achieves temporal layering for the tool-call event class
\(\mathcal{E}_{\text{tool}}\) in OpenClaw under three preconditions:

\textbf{(P1)} OpenClaw's plugin registry is read live on every dispatch
(the fix in PR-1 to issue \#5513).

\textbf{(P2)} The Kavach plugin registers on \texttt{before\_tool\_call}
during the platform's plugin-load phase, before any user-installable
extension can register tools.

\textbf{(P3)} The \texttt{before\_tool\_call} hook is invoked by the
runtime in \texttt{executeToolCalls()} (the fix in PR-1 to issue \#5943).

\begin{theorem}[Kavach's temporal coverage of tool calls]
\label{thm:coverage}
Under preconditions P1, P2, and P3, every tool call emitted by any agent
running in OpenClaw with the Kavach plugin installed is evaluated by the
Kavach parliament before its tool handler is invoked.
\end{theorem}

\paragraph{Proof sketch.} By P3, every tool call routes through
\texttt{executeToolCalls()}, which dispatches \texttt{before\_tool\_call}
before invoking any tool handler. By P1, dispatch reads the live plugin
registry; if the Kavach plugin is registered at dispatch time, its handler
is invoked. By P2, the Kavach plugin is registered in OpenClaw's
plugin-load phase, which precedes any agent run. There is no code path
that adds plugins after the agent run begins (P2's contract is enforced
by the platform). Hence at every tool-call dispatch, the Kavach handler
is in the dispatch list, and the tool handler runs only after the
parliament returns ALLOW or MODIFY. \(\square\)

The proof is mechanical, not deep. Its value lies in what it makes
explicit: temporal coverage rests on three plumbing properties of the
host platform. If any of P1, P2, P3 fails, the parliament still runs but
its verdict can be bypassed. The implementation work in this paper —
PR-1's two bug fixes plus the plugin's registration timing — is precisely
the work of establishing P1 and P3 (P2 holds by OpenClaw's load order).

\subsection{Asymmetric verdict combination}
\label{sec:asymmetric}

The speaker that combines minister verdicts uses a combination rule that has
not appeared by name in the prior guardrail literature; we define it here so
it can be cited precisely and distinguished from related but distinct prior
approaches.

\begin{definition}[Asymmetric verdict combination]
\label{def:asymmetric}
Let $E_1, \dots, E_k$ be enforcement points (ministers) producing verdicts in
$\{\textsc{Block}, \textsc{Escalate}, \textsc{Allow}\}$, and let $O$ be an
alignment oracle (COMPASS) producing a binary drift signal $\delta \in \{0,1\}$.
A verdict combiner is \emph{asymmetric} if:
\begin{enumerate}
  \item Any $E_i$ returning \textsc{Block} yields a final verdict of
        \textsc{Block}, regardless of the verdicts of all other $E_j$ and
        regardless of $\delta$. (\emph{Deny-first}: one signal suffices to block.)
  \item No score averaging or majority vote is computed across ministers.
  \item The oracle signal $\delta$ can \emph{escalate} a borderline
        (\textsc{Escalate}) verdict to a \textsc{Block} when combined with
        any minister returning \textsc{Escalate}, but cannot override a
        \textsc{Block} verdict from a minister.
\end{enumerate}
\end{definition}

\paragraph{Why not averaging?} The four ministers target structurally different
attack categories (code execution, credential theft, data exfiltration,
trajectory drift). A high cosine similarity from one minister is a meaningful
signal regardless of whether the other ministers return low scores — those
ministers are querying completely different corpora, so the absence of a match
in their domain is uninformative about the presence of a match in the triggered
minister's domain. Score averaging would dilute this signal in proportion to
the number of ministers that return low scores.

\paragraph{Relationship to prior work.} BARRED \cite{barred2026} uses the word
"asymmetric" to describe debate \emph{roles} during synthetic data generation,
not the verdict combination at inference time. ShieldAgent
\cite{chen2025shieldagent} combines probabilistic rule circuits via weighted
vote. RTBAS \cite{zhong2025rtbas} thresholds on an aggregate screener score.
None of these is asymmetric in the sense of Definition~\ref{def:asymmetric}.
Kavach's asymmetric combiner is closest to the $k$-out-of-$n$ circuit with
$k=1$ (any-of) from fault-tolerant systems theory, applied to security verdicts
rather than error signals.

\subsection{Spatial coverage and complementarity}

Within the parliament itself, Kavach is spatially layered with \(k=4\)
(four ministers + COMPASS as a fifth signal). The speaker's deny-first
combination rule means a single confident BLOCK from any minister
suffices to block, which is the strict spatial guarantee. We make no
claim that the four ministers' detection failure modes are independent;
measuring this is a stated future-work item (\S\ref{sec:future}) and a
legitimate threat to the spatial component of our claim.

The temporal and spatial axes are complementary, not substitutable. A
purely temporal defense (one enforcement point, always authoritative)
fails when that point misclassifies. A purely spatial defense fails when
extension code registers after initialization. Kavach combines both:
temporal coverage by virtue of the OpenClaw plugin contract (post PR-1),
spatial coverage by virtue of the parliament's four-minister structure.

This decomposition also clarifies what Kavach is \emph{not} a substitute
for. LlamaFirewall's PromptGuard 2 fires at message-input time, well
before any tool-call is resolved; Kavach's parliament fires at tool-call
time, after the model has chosen what to do. A deployment that ran both
would have spatial defense at \emph{two distinct enforcement boundaries}
along the temporal axis — input-time and tool-time — which we expect to
be a strictly stronger composition than either alone.

\subsection{Implications for the broader guardrail field}

The temporal--spatial distinction reframes a methodological question that
has been implicit in the literature: when comparing two guardrails on
the same benchmark, are we comparing two spatial detectors at the same
enforcement boundary, or two systems that enforce at different boundaries?
The numbers in InjecAgent and AgentDojo do not distinguish these cases,
because they replay attacks against the agent end-to-end and measure
whether the attack succeeded. A guardrail can succeed by spatial detection
at any boundary along the path.

We propose that future evaluations report not just the final ASR but also
\emph{at which enforcement boundary} the guardrail intervened, and
\emph{what time-step} the intervention occurred relative to component
registration. Without these axes, two guardrails with the same ASR on the
same benchmark may be doing structurally different work, and a deployment
choosing between them needs to know which.
